\documentclass[11pt]{article}

%%%%%%%%%%%%%%%%%%
% Basic packages %
%%%%%%%%%%%%%%%%%%

\usepackage[T1]{fontenc}
\usepackage[utf8]{inputenc}
\usepackage[british]{babel}
\usepackage[stretch = 25, shrink = 25, tracking=true, letterspace=30]{microtype}

\usepackage{geometry}
 \geometry{
 a4paper,
 %total={150mm,257mm},
 left=25mm,
 right=25mm,
 top=15mm,
 } % Page Layout
\usepackage[dvipsnames]{xcolor} % Colour
\usepackage{amsmath} % Maths

%%%%%%%%%%%%%%%%%%%
% Tables, figures %
%%%%%%%%%%%%%%%%%%%

\usepackage{float}
\usepackage{graphicx} % Imgages, figures
\graphicspath{{./figures/}} % Set figure path

\usepackage{booktabs} % Pretty tables

%%%%%%%%%%%%%
% Citations %
%%%%%%%%%%%%%

\usepackage[
    backend=biber,
    style=apa,
    sortlocale=de_DE,
    natbib=true,
    url=false, 
    doi=true,
    eprint=false
]{biblatex}

%\addbibresource{ASSIGNMENT_1.bib} %Imports bibliography file
\usepackage{csquotes} % used for block quotes

%%%%%%%%%
% Fonts %
%%%%%%%%%

%%% Helvetica:
%\usepackage[scaled=0.95]{helvet} % Imports Helvetica, which is similar to Arial
%\renewcommand{\rmdefault}{phv} % Use Helvetica

%%% Atkinson Hyperlegible, nerdy font 
%\usepackage[sfdefault]{atkinson} %% Option 'sfdefault' if the base
%% font of the document is to be sans serif.
%\usepackage[T1]{fontenc}

%%%%%%%%%%%%%%%%%%
% Header setting %
%%%%%%%%%%%%%%%%%%

\usepackage{titlesec} % pretty headers
\titleformat{\section}
  {\normalfont\large\bfseries\color{Black}} % Formatting options
  {\thesection}{1em}{}

%%%%%%%%%%%%%%%%%%
% Title settings %
%%%%%%%%%%%%%%%%%%

\title{\textcolor{NavyBlue}{\textbf{COMP527 - CA Assignment 2 - Preprocessing}}}
\author{\textbf{Pablo Ramos - 201885602}}

\setlength\parindent{10mm}

%=========%%%%%%%%%%%=========%
%=========%% BODY %%%=========%
%=========%%%%%%%%%%%=========%

\begin{document}
% Dark Mode - Comment the bottom 2 lines to get a normal white page
%\pagecolor{black}
%\color{white} % Set the default colour to white
%
%\pagenumbering{gobble} % Turns off page numbering
%
\maketitle
{\setlength\parindent{0mm} {\color{NavyBlue} \rule{\linewidth}{0.5mm}}} % Lines with colour options
\section{About the Data}

The data has 1760 sentences from various texts. A quick read of some of them shows that it's snippets from Shakespeare's Coriolanus, a New York noir novel, food-related writing and other miscellaneous or nonsensical phrases.

\section{Loading the Data}

The data was loaded as a Pandas data frame using the \texttt{pandas} library. The file is tab-separated, and \texttt{quoting=csv.QUOTE\_NONE} was set to prevent quote characters within sentences from being misinterpreted as delimiters.

\section{Cleaning the Data}

Text was cleaned using a combination of regex, spaCy, and NLTK. First, literal \texttt{\textbackslash n} sequences and non-alphabetic characters were removed with regex, and the text was lowercased. SpaCy's \texttt{en\_core\_web\_sm} model was then used to lemmatise each token, which reduces words to their root form (I.E., no declensions, grammatical cases, \&c.). 

The parser and named entity recognition components were disabled to speed up the script (not enough compute).

NLTK's English stopword list was applied to discard common words with little semantic value. Tokens shorter than three characters were also removed. 

Sentences left empty after cleaning were replaced with the placeholder "empty" to preserve row alignment with the original dataset.

\section{Lemmatisation}

Lemmatisation reduces vocabulary size and consolidates related word forms, improving the quality of sentence embeddings. Stopword removal ensures that high-frequency structural words do not dominate the representation. 

These steps were picked because they produce cleaner input for the SBERT model, which will be covered in report 2.

%\begin{displayquote}
%    This is some text from a book.
%\end{displayquote}

%%%%%%%%%%%%%%%%%%%%%%
% Implicit citations %
%%%%%%%%%%%%%%%%%%%%%%

%\nocite{Hunter:2007}
% \nocite{Waskom2021}


 %\printbibliography

\end{document}